% !TeX root = ../qm_chap3.tex

\begin{frame}{教材3.3　测得本征值的概率由投影决定}
  \iqbfontsizemode{small}
    \iqbsectiontitle{测量结果的量子化}
    \begin{iqbitemize}
    \item 对离散谱，结果是$\hat A$的某个本征值$a_n$。初态若在该本征子空间，测量必得$a_n$。
    \item \textbf{玻恩规则}（Born rule）：
    \subitem{对归一化态$|\psi\rangle$，非简并结果$a_n$的概率等于对应展开系数的模平方。}
    \end{iqbitemize}
      \[
        P_n = |\langle\phi_n|\psi\rangle|^2 = |c_n|^2
      \]
    \begin{iqbitemize}
    \item 其中$c_n$是展开系数：$|\psi\rangle = \sum_n c_n|\phi_n\rangle$
    \end{iqbitemize}
    \iqbsmallsep

    \iqbsectiontitle{态坍缩（wave function collapse）}
    \begin{iqbitemize}
    \item 测量前：系统处于叠加态$|\psi\rangle = \sum_n c_n|\phi_n\rangle$
    \item 若$a_n$简并，应把投影到该子空间的各正交分量概率相加。
    \item 理想投影测量后，状态是相应子空间内的归一化投影，无时间演化时立即重测得到同一结果。
    \end{iqbitemize}
  
\end{frame}

\begin{frame}{教材3.3　初态决定哪些箱中能级会被测到}
  \iqblayouttwo{
    \iqbsectiontitle{抛物线初态是多个定态的叠加}
    \begin{iqbitemize}
      \item 在$0\le x\le l$内，教材的初态为$\Psi(x,0)=\sqrt{30/l^5}\,x(l-x)$，箱外为零。
      \item 能量测量的候选值为$E_n=n^2h^2/(8ml^2)$。初态与箱中心镜面对称，只和同为偶宇称的$\psi_n$有非零重叠，即$n$为奇数。
      \item 将初态投影到$\psi_n=\sqrt{2/l}\sin(n\pi x/l)$，得到$c_n$。各允许能级以$|c_n|^2$出现。
    \end{iqbitemize}
  }{
    \iqbsectiontitle{第三个定态的能量确定}
    \begin{iqbitemize}
      \item 若初态为$\sqrt{2/l}\sin(3\pi x/l)$，它就是归一化的$\psi_3$，展开中只有$c_3=1$。
      \[
        E_3=\dfrac{9h^2}{8ml^2},\quad P(E_3)=1 .
      \]
      \item 每次实验若都制备相同初态，第一种制备产生能量的统计分布，第二种制备每次给出同一个能量。
    \end{iqbitemize}
  }
\end{frame}

\begin{frame}{可观测量的测量（2/2）}
  \iqbfontsizemode{small}
  \iqblayouttwo{
    \iqbsectiontitle{期望值与不确定性}
    \begin{iqbitemize}
    \item \textbf{期望值}（expectation value）：
    \subitem{大量相同系统测量的统计平均值}
    \end{iqbitemize}
      \[
        \langle A \rangle = \sum_n a_n P_n = \langle\psi|\hat{A}|\psi\rangle
      \]
    \begin{iqbitemize}
    \item 狄拉克符号表示最简洁
    \item \textbf{不确定性}（uncertainty）：
    \end{iqbitemize}
      \[
        (\Delta A)^2 = \langle A^2 \rangle - \langle A \rangle^2
      \]
    \begin{iqbitemize}
    \item 测量结果相对于期望值的统计分散度
    \item 标准差：$\Delta A = \sqrt{\langle A^2 \rangle - \langle A \rangle^2}$
    \end{iqbitemize}
  }{
    \iqbsectiontitle{测量理论的含义}
    \begin{iqbitemize}
    \item \textbf{概率性}：
    \subitem{量子力学本质上是概率理论}
    \subitem{不同于经典力学的决定论}
    \item \textbf{测量干扰}：
    \subitem{一般叠加态经投影测量会发生改变}
    \subitem{若已处在对应的本征子空间，理想测量保持该状态}
    \item \textbf{客观实在性}：
    \subitem{测量前物理量通常没有确定值}
    \subitem{只有测量后才"坍缩"到确定状态}
    \end{iqbitemize}
    \iqbtinysep

    \iqbgreenbox{哲学意义}{量子测量挑战了经典的客观实在概念，引发关于观察者角色的深刻讨论}
  }
\end{frame}

\begin{frame}{三维薛定谔方程（1/2）}
  \iqbfontsizemode{small}
  \iqblayouttwo{
    \iqbsectiontitle{薛定谔方程的三维形式}
    \begin{iqbitemize}
    \item \textbf{含时薛定谔方程}：
    \end{iqbitemize}
      \[
        i\hbar\dfrac{\partial}{\partial t}\Psi(x,y,z,t) = \hat{H}\Psi(x,y,z,t)
      \]
    \begin{iqbitemize}
    \item \textbf{定态薛定谔方程}（time-independent）：
    \end{iqbitemize}
      \[
        \hat{H}\psi(x,y,z) = E\psi(x,y,z)
      \]
    \begin{iqbitemize}
    \item 波函数依赖三个空间坐标：$\Psi = \Psi(x,y,z,t)$
    \item 当$V=V(x,y,z)$与时间无关时，定态的时间因子为$e^{-iEt/\hbar}$。
    \end{iqbitemize}
  }{
    \iqbsectiontitle{三维哈密顿算符}
    \begin{iqbitemize}
    \item 单粒子的哈密顿算符：
    \end{iqbitemize}
      \[
        \hat{H} = \hat{T} + \hat{V} = -\dfrac{\hbar^2}{2m}\nabla^2 + V(x,y,z)
      \]
    \begin{iqbitemize}
    \item $\hat{T}$：动能算符（kinetic energy operator）
    \item $\nabla^2$：拉普拉斯算符（Laplacian operator）
    \end{iqbitemize}
      \[
        \nabla^2 = \dfrac{\partial^2}{\partial x^2} + \dfrac{\partial^2}{\partial y^2} + \dfrac{\partial^2}{\partial z^2}
      \]
    \begin{iqbitemize}
    \item 描述空间中的二阶微分，是梯度的散度
    \end{iqbitemize}
  }
\end{frame}

\begin{frame}{三维薛定谔方程（2/2）}
  \iqbfontsizemode{small}
    \iqbsectiontitle{球坐标系中的表达}
    \begin{iqbitemize}
    \item 对于中心力场（如氢原子），使用球坐标$(r,\theta,\phi)$更方便
    \item 球坐标拉普拉斯算符：
    \end{iqbitemize}
      \[
        \nabla^2 = \dfrac{1}{r^2}\dfrac{\partial}{\partial r}\left(r^2\dfrac{\partial}{\partial r}\right) + \dfrac{1}{r^2\sin\theta}\dfrac{\partial}{\partial\theta}\left(\sin\theta\dfrac{\partial}{\partial\theta}\right) + \dfrac{1}{r^2\sin^2\theta}\dfrac{\partial^2}{\partial\phi^2}
      \]
    \begin{iqbitemize}
    \item 分离变量：$\psi(r,\theta,\phi) = R(r)\Theta(\theta)\Phi(\phi)$
    \end{iqbitemize}
    \iqbsmallsep

    \iqbsectiontitle{多粒子系统}
    \begin{iqbitemize}
    \item $N$个粒子的薛定谔方程：
    \end{iqbitemize}
      \[
        i\hbar\dfrac{\partial}{\partial t}\Psi(\mathbf{r}_1,\ldots,\mathbf{r}_N,t) = \left[-\sum_{i=1}^{N}\dfrac{\hbar^2}{2m_i}\nabla_i^2 + V(\mathbf{r}_1,\ldots,\mathbf{r}_N,t)\right]\Psi
      \]
    \begin{iqbitemize}
    \item 波函数依赖$3N$个空间坐标
    \item 例如：氢分子（2个电子）依赖6个坐标
    \item 维数灾难：计算复杂度随粒子数指数增长
    \end{iqbitemize}
    \iqbmicrosep

    \iqbgreenbox{实际意义}{三维薛定谔方程是描述原子、分子和量子化学的基本方程}
  

\end{frame}

\begin{frame}{教材3.4　多粒子势能可以依赖粒子间距离}
  \iqblayouttwo{
    \iqbsectiontitle{把相互作用写成坐标的函数}
    \begin{iqbitemize}
      \item 对两个相互作用的粒子，波函数$\psi(\mathbf r_1,\mathbf r_2)$依赖六个空间坐标。
      \item 教材例子取$V=C/r_{12}$，其中$r_{12}=|\mathbf r_1-\mathbf r_2|$，由三个坐标差的平方和开根得到。
      \item $V$以乘法形式作用于波函数，通常会耦合两个粒子的运动。
    \end{iqbitemize}
  }{
    \iqbsectiontitle{每个动能项只对自身坐标求导}
    \[
      \left[-\dfrac{\hbar^2}{2m_1}\nabla_1^2
            -\dfrac{\hbar^2}{2m_2}\nabla_2^2
            +\dfrac{C}{r_{12}}\right]\psi=E\psi .
    \]
    \begin{iqbitemize}
      \item $\nabla_i^2$只对$(x_i,y_i,z_i)$求二阶偏导，另一粒子的坐标在该项中保持固定。
      \item 电子与质子之间的库仑吸引对应$C<0$。改用质心和相对坐标可进一步分离运动。
    \end{iqbitemize}
  }
\end{frame}

\begin{frame}{教材3.4　概率密度须乘以体积元}
  \iqblayouttwo{
    \iqbsectiontitle{一个粒子的三维位置概率}
    \[
      \mathrm{d}P=|\Psi(\mathbf r,t)|^2\,\mathrm{d}^3r,\qquad
      \mathrm{d}^3r=\mathrm{d}x\,\mathrm{d}y\,\mathrm{d}z .
    \]
    \iqbfigcap[width=\linewidth,height=0.39\textheight]{Figures/fig3_1.png}{图3.1　对体积元内的概率密度积分，得到在该小区域内发现粒子的概率。}
  }{
    \iqbsectiontitle{多个粒子对应组态空间中的联合概率}
    \begin{iqbitemize}
      \item 对$N$个粒子，$3N$维体积元是各粒子体积元的乘积：
      \[
        \mathrm{d}\tau=\prod_{i=1}^N\mathrm{d}^3r_i .
      \]
      \item $|\Psi|^2\mathrm{d}\tau$给出所有粒子同时落入各自小区域的联合概率。
      \item 在全部坐标的取值范围上积分，总概率为1。
      \[
        \int|\Psi(\mathbf r_1,\ldots,\mathbf r_N,t)|^2\,\mathrm{d}\tau=1 .
      \]
      \item 对定态，时间因子的模平方为1，可直接用空间波函数$\psi$归一化。
    \end{iqbitemize}
  }
\end{frame}

\begin{frame}{教材3.5　三维无限深势阱的波函数在六面墙上为零}

  \iqblayouttwo{
    \iqbsectiontitle{物理模型}
    \begin{iqbitemize}
      \item 粒子限制在长方体箱子中
      \item 边长：$a$, $b$, $c$
      \item 箱内$V=0$，箱外$V=\infty$
      \item 边界条件：
      \[
        \psi(0,y,z) = \psi(a,y,z) = 0
      \]
      \[
        \psi(x,0,z) = \psi(x,b,z) = 0
      \]
      \[
        \psi(x,y,0) = \psi(x,y,c) = 0
      \]
    \end{iqbitemize}
  }{
    \iqbsectiontitle{长方体势阱}
    \iqbfigcap[width=\linewidth,height=0.5\textheight]{Figures/fig3_2.png}{图3.2　箱内势能$V=0$，箱外波函数为零。由连续性，波函数在每面墙上也为零。}
  }
\end{frame}

\begin{frame}{教材3.5　三项分别只依赖一个坐标，故各自为常数}
  \iqblayouttwo{
    \iqbsectiontitle{代入乘积形式}
    \begin{iqbitemize}
      \item 在箱内令$\psi=XYZ$，其中$X$只依赖$x$、$Y$只依赖$y$、$Z$只依赖$z$。
      \item 代入箱内薛定谔方程，在乘积非零处同除以$XYZ$。
      \[
        -\dfrac{\hbar^2}{2m}
        \left(\dfrac{X''}{X}+\dfrac{Y''}{Y}+\dfrac{Z''}{Z}\right)=E .
      \]
      \item 固定$y,z$并改变$x$，前一项也须保持不变。对$y,z$同理，定义三个分离常数$E_x,E_y,E_z$。
    \end{iqbitemize}
  }{
    \iqbsectiontitle{每个分量满足同一个一维箱问题}
    \[
      X''+\dfrac{2mE_x}{\hbar^2}X=0,\qquad X(0)=X(a)=0 .
    \]
    \begin{iqbitemize}
      \item 对$y,z$分别用边长$b,c$和常数$E_y,E_z$。
      \item 三个一维能量相加得到$E=E_x+E_y+E_z$。波函数的归一化因子等于三个一维因子的乘积。
      \item 三个方向的乘积本征函数构成一组基底。简并时可将同一能级的乘积解作线性组合。
    \end{iqbitemize}
  }
\end{frame}

\begin{frame}{三维箱中的粒子：数学求解}
  \iqblayouttwo{
    \iqbsectiontitle{分离变量求解}
    \begin{iqbitemize}
      \item 设$\psi(x,y,z) = X(x)Y(y)Z(z)$
      \item 分解为三个独立的一维问题
      \item \textbf{本征函数}：
      \begin{align*}
        \psi_{n_x,n_y,n_z}(x,y,z) &= \sqrt{\dfrac{8}{abc}} \sin\left(\dfrac{n_x\pi x}{a}\right) \\
        &\quad \times \sin\left(\dfrac{n_y\pi y}{b}\right) \sin\left(\dfrac{n_z\pi z}{c}\right)
      \end{align*}

      \item \textbf{能量本征值}：
      \begin{align*}
        E_{n_x,n_y,n_z} &= \dfrac{h^2}{8m}\left(\dfrac{n_x^2}{a^2} + \dfrac{n_y^2}{b^2} + \dfrac{n_z^2}{c^2}\right) \\
        n_x, n_y, n_z &= 1, 2, 3, \ldots
      \end{align*}
    \end{iqbitemize}

    \iqbtinysep

    \iqborangebox{关键结论}{三维能级依赖于三个量子数的平方和，体现各向异性}
  }{
    \iqbsectiontitle{物理意义}
    \begin{iqbitemize}
      \item \textbf{量子数的独立性}
      \begin{iqbitemize}
        \item $n_x$, $n_y$, $n_z$各自独立取值
        \item 总能量是三个方向能量之和
      \end{iqbitemize}

      \item \textbf{各向异性箱}
      \begin{iqbitemize}
        \item 三边互不相等时，交换量子数通常改变能量
        \item 某些边长比例仍允许不同量子数组合给出相同能量，称为偶然简并
      \end{iqbitemize}

      \item \textbf{与经典对比}
      \begin{iqbitemize}
        \item 经典：能量连续可变
        \item 量子：能量量子化，由三个整数决定
      \end{iqbitemize}
    \end{iqbitemize}
  }
\end{frame}

\begin{frame}{立方箱的能量只依赖三个量子数的平方和}
  \iqbfontsizemode{small}

  \iqbsectiontitle{立方箱条件}
  \begin{iqbitemize}
    \item \textbf{高度对称}：$a = b = c$

    \item \textbf{能量公式}：
    \[
      \begin{aligned}
        E_{n_x,n_y,n_z} &= \dfrac{h^2}{8ma^2} \\
        &\quad \times (n_x^2 + n_y^2 + n_z^2)
      \end{aligned}
    \]

    \item \textbf{简并实例}：
    \begin{iqbitemize}
      \item $E_{1,1,1}$：无简并（基态）
      \item $\begin{aligned}[t]
        E_{1,1,2} \\
        = E_{1,2,1} \\
        = E_{2,1,1}
      \end{aligned}$（三重简并）
      \item $\begin{aligned}[t]
        E_{1,2,2} \\
        = E_{2,1,2} \\
        = E_{2,2,1}
      \end{aligned}$（三重简并）
    \end{iqbitemize}

    \item \textbf{图示说明}：基态$E_{1,1,1}$无简并，激发态如$E_{1,1,2}$、$E_{1,2,2}$出现多重简并

    \item \textbf{物理原因}：
    立方箱允许置换坐标轴，量子数的不同排列可以给出同一能量
  \end{iqbitemize}
\end{frame}

\begin{frame}{立方箱的前两个激发能级均为三重简并}
  \iqbfontsizemode{small}
  \iqblayouttwo{
    \iqbsectiontitle{图示要点}
    \begin{iqbitemize}
      \item 基态无简并，激发态出现多重简并
      \item 相同高度能级对应不同的量子数组合
      \item 图中两个激发能级均为三重简并。简并度由对应的量子数组合数决定
    \end{iqbitemize}
  }{
    \iqbsectiontitle{能级图示}
    \iqbfigcap[width=\linewidth,height=0.55\textheight]{Figures/fig3_3.png}{图3.3　立方箱中粒子的最低三个能级。横线下的数字依次表示$n_x$、$n_y$和$n_z$。}
  }
\end{frame}

\begin{frame}{教材3.6　同一能级的非零叠加仍具有确定能量}
  \iqblayouttwo{
    \iqbsectiontitle{用线性证明简并子空间的性质}
    \begin{iqbitemize}
      \item 若每个$\psi_j$都满足$\hat H\psi_j=E\psi_j$，对非零叠加$\phi=\sum_j c_j\psi_j$有
      \[
        \hat H\phi=\sum_j c_j\hat H\psi_j
                  =E\sum_j c_j\psi_j=E\phi .
      \]
      \item 立方箱中$\psi_{211}$、$\psi_{121}$和$\psi_{112}$的非零叠加仍有能量$6h^2/(8ma^2)$。
      \item 两个不同能级的本征函数都以非零系数出现时，不能提出同一个能量因子，叠加态能量不确定。
    \end{iqbitemize}
  }{
    \iqbsectiontitle{简并度数的是线性无关的态}
    \begin{iqbitemize}
      \item 若$\sum_j c_j\psi_j=0$迫使所有$c_j=0$，这些波函数线性无关。倍乘同一个波函数不会增加简并度。
      \item 一维自由粒子在同一个正能量$E=p^2/(2m)$下有两个独立方向，对应$e^{ipx/\hbar}$和$e^{-ipx/\hbar}$。
      \item 可在简并子空间内重新选择正交基底，能量保持不变。
    \end{iqbitemize}
  }
\end{frame}

\begin{frame}{例题3.3：验证对易关系$[\hat{x}, \hat{p}_x] = i\hbar$（1/2）}
  \iqbfontsizemode{small}
  \iqbsectiontitle{最基本的量子对易关系}

  \textbf{问题}：计算对易子$[\hat{x}, \hat{p}_x]$，其中$\hat{p}_x = -i\hbar\dfrac{\mathrm{d}}{\mathrm{d}x}$

  \begin{iqbitemize}
    \item \textbf{对易子定义}：
    \[
      [\hat{x}, \hat{p}_x] = \hat{x}\hat{p}_x - \hat{p}_x\hat{x}
    \]

    \item \textbf{作用于任意函数$f(x)$}：
    \begin{align*}
      [\hat{x}, \hat{p}_x]f(x) &= \hat{x}\left(-i\hbar\dfrac{\mathrm{d}}{\mathrm{d}x}\right)f(x) - \left(-i\hbar\dfrac{\mathrm{d}}{\mathrm{d}x}\right)\hat{x}f(x) \\
      &= -i\hbar x\dfrac{\mathrm{d}f}{\mathrm{d}x} - \left[-i\hbar\dfrac{\mathrm{d}}{\mathrm{d}x}(xf)\right] \\
      &= -i\hbar x\dfrac{\mathrm{d}f}{\mathrm{d}x} + i\hbar\dfrac{\mathrm{d}}{\mathrm{d}x}(xf)
    \end{align*}
  \end{iqbitemize}

  \iqbgreenbox{关键步骤}{先展开对易子，再作用于任意函数}
\end{frame}

\begin{frame}{例题3.3：验证对易关系$[\hat{x}, \hat{p}_x] = i\hbar$（2/2）}
  \iqbfontsizemode{small}
  \iqbsectiontitle{乘积法则的应用}

  \begin{iqbitemize}
    \item \textbf{使用乘积法则}：$\dfrac{\mathrm{d}}{\mathrm{d}x}(xf) = f + x\dfrac{\mathrm{d}f}{\mathrm{d}x}$
    \begin{align*}
      &= -i\hbar x\dfrac{\mathrm{d}f}{\mathrm{d}x} + i\hbar\left(f + x\dfrac{\mathrm{d}f}{\mathrm{d}x}\right) \\
      &= -i\hbar x\dfrac{\mathrm{d}f}{\mathrm{d}x} + i\hbar f + i\hbar x\dfrac{\mathrm{d}f}{\mathrm{d}x} \\
      &= i\hbar f(x)
    \end{align*}

    \item \textbf{结论}：
    \[
      [\hat{x}, \hat{p}_x]f(x) = i\hbar f(x) \quad \Rightarrow \quad [\hat{x}, \hat{p}_x] = i\hbar
    \]
  \end{iqbitemize}

  \iqborangebox{核心意义}{此对易关系是海森堡不确定性原理的数学根源}
\end{frame}

\begin{frame}{例题3.4：对易子的计算技巧}
  \iqbfontsizemode{small}
  \iqbsectiontitle{利用对易关系简化计算}

  \textbf{问题}：计算$[\hat{x}, \hat{p}_x^2]$

  \begin{iqbitemize}
    \item \textbf{方法一：直接计算}
    \begin{align*}
      [\hat{x}, \hat{p}_x^2] &= \hat{x}\hat{p}_x^2 - \hat{p}_x^2\hat{x} \\
      &= \hat{x}\hat{p}_x\hat{p}_x - \hat{p}_x\hat{p}_x\hat{x}
    \end{align*}

    \item \textbf{方法二：利用对易子恒等式}
    \[
      [\hat{A}, \hat{B}\hat{C}] = [\hat{A}, \hat{B}]\hat{C} + \hat{B}[\hat{A}, \hat{C}]
    \]

    \item \textbf{应用恒等式}：
    \begin{align*}
      [\hat{x}, \hat{p}_x^2] &= [\hat{x}, \hat{p}_x]\hat{p}_x + \hat{p}_x[\hat{x}, \hat{p}_x] \\
      &= i\hbar\hat{p}_x + \hat{p}_x(i\hbar) \\
      &= i\hbar\hat{p}_x + i\hbar\hat{p}_x \\
      &= 2i\hbar\hat{p}_x
    \end{align*}

    \item \textbf{推广}：
    \[
      [\hat{x}, \hat{p}_x^n] = ni\hbar\hat{p}_x^{n-1}
    \]
    可以用数学归纳法证明
  \end{iqbitemize}

  \iqbgreenbox{技巧}{利用对易子恒等式可大大简化计算}
\end{frame}
